\section{Round-eleven positive closure: dynamic cohomology and covariance forms}
\label{sec:r11-b3}

The covariance is constructed first from finite-volume cumulants.  No
curvature is attributed to the linear balance multiplier, and no inverse is
placed on the closure of the covariance range.

\subsection{Regular biased paths and hypocoercivity}
A regular exposed path is uniformly comparable to a local Maxwellian, has
bounded weighted derivatives, and has collision ratio
$0<q_-\le q\le q_+$.  These conditions are verified on compact bounded
source charts and form a rate-dense class by B2 regularization.

\begin{theorem}[Inhomogeneous short-time energy estimate]
\label{thm:r11-b3-energy}
The linearized biased Boltzmann balance around a regular exposed path is
well-posed in the weighted tangent space.  Its microscopic part is controlled
by comparison of the symmetric collision form with the Maxwellian Dirichlet
form, and its spatially dependent macroscopic fields are controlled by a
Fourier hypocoercive estimate.  Consequently, for sufficiently short fixed
$T$, the constrained second variation is a closed coercive form modulo the
five local conservation fields and declared endpoint constraints.
\end{theorem}

\begin{proof}
Maxwellian comparability and the bounds on $q$ compare the local collision
Dirichlet form with the equilibrium form on the orthogonal complement of
invariants.  Fourier-transform in $x$; for every spatial frequency combine
the transport skew term with the microscopic coercivity by the standard
compensating cross term.  The zero frequency is fixed by initial constraints,
while nonzero frequencies gain a uniform short-time estimate.  The possibly
indefinite perspective term is bounded by $CT\|q-1\|$ times the tangent norm
and is absorbed after choosing the declared short-time chart.
\end{proof}

\subsection{The exact dynamic gauge}
Include endpoint cotangents explicitly.  Let $\mathfrak G r$ add the weak
balance test $r$, subtract $\Delta r$ from the collision source, and add its
initial and terminal traces to the endpoint cotangents.  Let temporal
coboundaries be the closure of $U-U\circ\Theta_t$ in the source norm.

\begin{theorem}[Representation and dynamic cohomology quotient]
\label{thm:r11-b3-gauge}
The annihilator of every balanced tangent is the closed sum of the exact
balance gauge $\operatorname{Ran}\mathfrak G$, collision invariants, and
temporal coboundaries.  The range is closed and the quotient is Hausdorff.
\end{theorem}

\begin{proof}
Integration by parts in the weak balance, including the two endpoint traces,
shows that every gauge vector annihilates the tangent.  The backward
hypocoercive estimate dual to Theorem~\ref{thm:r11-b3-energy} gives a bounded
right inverse after the conservation and endpoint normalizations, hence
closed range.  If a source annihilates all balanced tangents, solve the
adjoint balance equation and subtract its gauge representative.  The
remaining zero-variance additive source is a temporal coboundary by the
regular-phase Livsic theorem obtained from the B2 local source pressure.
\end{proof}

\subsection{Covariance first and the Cameron--Martin form}

\begin{theorem}[Finite-volume covariance and process limit]
\label{thm:r11-b3-clt}
The exact-centered density/contact fluctuation field has finite-dimensional
cumulants converging to a positive covariance form $\Sigma$.  Time-localized
B2 source derivatives give, for every deterministic interval,
\[
 \mathbb E|Z_t-Z_s|^4\le C|t-s|^2,
\]
including all density--contact cross terms and contact diagonals.  Mitoma and
Kolmogorov therefore yield a continuous Gaussian process in the weighted
nuclear dual.
\end{theorem}

\begin{proof}
Differentiate the normally convergent B2 cumulant on disjoint time-source
blocks.  Connected trees meeting $k$ prescribed blocks carry the volume of
their ordered time hull, which gives the interval factors.  Third and higher
central-limit cumulants vanish with the power of $\mu_\varepsilon$; the
second cumulant converges.  The fourth-moment estimate and compact containment
prove tightness, while finite-dimensional convergence identifies the limit.
\end{proof}

\begin{theorem}[Second epi-derivative on the correct domain]
\label{thm:r11-b3-mosco}
The second epi-derivative of the constrained rate at a regular exposed path is
the closed extended quadratic form
\[
 \mathfrak q(z)=
 \begin{cases}
  \|\Sigma^{-1/2}z\|^2,&z\in\operatorname{Ran}\Sigma^{1/2},\\
  +\infty,&\text{otherwise}.
 \end{cases}
\]
Its nullspace is exactly the dynamic cohomology quotient of
Theorem~\ref{thm:r11-b3-gauge}.  On the rate-dense regular class this form
agrees with the closure of the correctly differentiated B2 action.
\end{theorem}

\begin{proof}
On every finite source projection, analytic convex duality identifies the
second epi-derivative with the inverse covariance matrix after the exact B1
Schur complement.  The projective forms are equicoercive in the energy norm of
Theorem~\ref{thm:r11-b3-energy}; diagonal recovery gives Mosco convergence.
The spectral theorem for positive operators identifies the limiting form
domain as $\operatorname{Ran}\Sigma^{1/2}$, not the closure of
$\operatorname{Ran}\Sigma$.  Direct differentiation of the relative-entropy
perspective on balanced tangents and the energy estimate identifies the same
closed form.
\end{proof}

\begin{theorem}[Round-eleven B3 closure]
\label{thm:r11-b3-main}
Regular hard-sphere phases possess a well-posed inhomogeneous linearized
balance, an exact representation-plus-temporal cohomology quotient, a
finite-cumulant Gaussian process, and a closed Cameron--Martin inverse form on
$\operatorname{Ran}\Sigma^{1/2}$.
\end{theorem}

\begin{proof}
Combine Theorems~\ref{thm:r11-b3-energy},
\ref{thm:r11-b3-gauge}, \ref{thm:r11-b3-clt}, and
\ref{thm:r11-b3-mosco}.
\end{proof}
